% MA 214 Discussion 10 -- covers Lectures 14 and 15.
%
% Student handout. Page 1 is the concept review; the question set runs from page 2.
% Compile from inside this folder:  pdflatex Discussion10.tex

\documentclass[11pt]{article}
\usepackage[margin=1in]{geometry}
\usepackage{parskip}
\usepackage{fancyhdr}
\usepackage{array}
\usepackage{booktabs}
\usepackage{enumitem}
\usepackage{amsmath}
\usepackage{tabularx}
\usepackage{listings}

\pagestyle{fancy}
\fancyhf{}
\cfoot{\thepage}
\rhead{Discussion 10}
\lhead{MA 214: Applied Statistics}
\renewcommand{\headrulewidth}{.4mm}

\newcommand{\blankline}[1]{\underline{\hspace{#1}}}
\newcolumntype{L}{>{\raggedright\arraybackslash}X}

\lstset{
  basicstyle=\ttfamily\small,
  columns=fullflexible,
  frame=single,
  framerule=0.4pt,
  breaklines=true,
  showstringspaces=false,
  aboveskip=8pt,
  belowskip=8pt
}

% Section header: bold title over a hairline rule.
\newcommand{\parttitle}[1]{%
  \par\medskip\noindent
  {\large\bfseries #1}\par
  \vspace{-.5em}\noindent\rule{\linewidth}{0.4pt}\par\smallskip}

% Writing space.
\newcommand{\answerspace}[1]{\vspace{#1}}

% Flush-left question list: the label runs in at the margin and the body wraps
% back to the margin, so nothing is pushed far to the right.
\newlist{questions}{enumerate}{1}
\setlist[questions]{label=\textbf{Question \arabic*.}, wide=0pt, itemsep=1.4em, topsep=.8em}

\newlist{parts}{enumerate}{1}
\setlist[parts]{label=\textbf{(\Alph*)}, leftmargin=1.6em, labelsep=.45em,
  itemsep=.7em, topsep=.5em}

\begin{document}

\noindent\textbf{Name:}\blankline{2.3in}\hfill\textbf{BUID:}\blankline{1.6in}

\begin{center}
  {\Large\bfseries Discussion 10: Interactions, Two Intervals, and the Cliff}
\end{center}

\noindent\textbf{Goals:} \textbf{(1)} read an interaction coefficient as a change in a slope,
\textbf{(2)} judge whether an extra term earned its place, \textbf{(3)} tell a confidence interval
from a prediction interval and name the three sources of doubt, \textbf{(4)} explain why
cross-validation sees overfitting that $R^2$ and AIC do not.

\parttitle{Part 1: Concept Review}

{\small

\noindent Everything in Part 2 can be answered from this page. Keep it in front of you.

\begin{enumerate}[label=\textbf{\arabic*.}, wide=0pt, itemsep=.4em, topsep=.5em]

\item \textbf{Every row is a conditional claim.} In a multiple regression each coefficient is the
predicted change in $y$ per one unit of that predictor, \emph{holding the others fixed}. Its
$t = \hat\beta/\text{SE}$ has df $= n - p - 1$, and its interval is
$\hat\beta \pm t^{\ast}\text{SE}$.

\item \textbf{An interaction is a slope that depends on something else.} In
$\hat{y} = \hat\beta_0 + \hat\beta_1 x + \hat\beta_2 g + \hat\beta_3 xg$ with $g$ an indicator,
\[
\text{slope in } x \text{ when } g=0 \;=\; \hat\beta_1,
\qquad
\text{slope in } x \text{ when } g=1 \;=\; \hat\beta_1 + \hat\beta_3 .
\]
So $\hat\beta_3$ is a \emph{difference of slopes}, not a difference of levels. And once an
interaction is present, $\hat\beta_2$ is only the group gap at $x = 0$; if $x = 0$ is outside the
data, that coefficient is a nonsense point and its $p$-value says nothing about whether the group
matters.

\item \textbf{Did the term earn its place?} Compare adjusted $R^2$ or AIC between the models, and
look at whether the interaction's own interval excludes $0$. A larger $R^2$ alone never settles it.

\item \textbf{Two intervals, two questions.} At a chosen $x^{\ast}$ both are centered on the same
$\hat{y}$:
\[
\textbf{CI for the mean: } \hat{y} \pm t^{\ast}\,\text{SE}_{\text{fit}},
\qquad
\textbf{PI for one new case: } \hat{y} \pm t^{\ast}\sqrt{\hat\sigma^2 + \text{SE}_{\text{fit}}^2}.
\]
The CI answers ``where is the average for branches like this?'' The PI answers ``where will
\emph{this one} land?'' The PI is always wider.

\item \textbf{Three sources of doubt, and what infinite data buys.} A prediction carries
(1) irreducible scatter $\hat\sigma^2$, (2) uncertainty in the fitted level, and (3) uncertainty in
the slope, which grows with distance from $\bar{x}$. More data shrinks (2) and (3) toward zero but
leaves (1) untouched, so the CI can shrink to nothing while the PI cannot get narrower than about
$\pm t^{\ast}\hat\sigma$.

\item \textbf{Overfitting and cross-validation.} $R^2$ can only rise as terms are added, so it
always prefers the biggest model. Adjusted $R^2$ and AIC charge for parameters and do better, but
both are computed on the \emph{same} data that fitted the model. In $k$-fold cross-validation each
observation is predicted by a model that never saw it, which is why CV can detect a model that fits
better and predicts worse. You cannot grade a model on the data that trained it.

\item \textbf{Four traps.} Reading an interaction coefficient as a group difference. Concluding a
group does not matter from the $p$-value of its main effect in an interaction model. Reporting a CI
when a single new case was asked about. Choosing among models of different size by $R^2$.

\end{enumerate}

}

\newpage

\parttitle{Part 2: Question Set}

\noindent\textbf{Scenario.} For $n = 90$ branches, annual visits (thousands) are modelled from
weekly open \texttt{hours} (range $20$ to $70$) and \texttt{type} (Branch or Central, with Branch as
the reference). Fitting \texttt{visits} $\sim$ \texttt{hours} $\ast$ \texttt{type}:

\begin{center}
\renewcommand{\arraystretch}{1.2}
\begin{tabular}{lcccc}
\toprule
 & \textbf{Estimate} & \textbf{Std. Error} & \textbf{$t$} & \textbf{$p$-value} \\
\midrule
(Intercept) & $17.036$ & 3.579 & $4.76$ & $<0.001$ \\
\texttt{hours} & $1.0303$ & 0.0782 & $13.18$ & $<0.001$ \\
\texttt{type}Central & $-3.219$ & 7.427 & $-0.43$ & $0.666$ \\
\texttt{hours:type}Central & $0.9588$ & 0.1535 & $6.25$ & $<0.001$ \\
\bottomrule
\end{tabular}
\end{center}

\noindent Also: $\hat\sigma = 9.297$, df $= 86$, $t^{\ast}_{86} = 1.9879$, adjusted
$R^2 = 0.907$. The main-effects model without the interaction has adjusted $R^2 = 0.867$ and
AIC $694.33$; this model has AIC $662.66$.

\begin{questions}

%%%%%%%%%%%%%%%%%%%%%%%%%%%%%%%%%%%%%%
\item \textbf{Two slopes, not one.}

\begin{parts}
\item Write the fitted equation for a Branch branch and for a Central branch, then give the slope
in \texttt{hours} for each.
\answerspace{4.5em}

\item Interpret $0.9588$ in context. Your sentence must make clear that it is a difference of
\emph{slopes}.
\answerspace{4em}

\item Compute the 95\% interval for the interaction coefficient. Does it exclude $0$?
\answerspace{3.5em}

\item The \texttt{type}Central coefficient has $p = 0.666$. A classmate concludes that branch type
does not matter. Given your answer to (A), explain why that is badly wrong, and say what
$-3.219$ actually refers to.
\answerspace{4.5em}
\end{parts}

%%%%%%%%%%%%%%%%%%%%%%%%%%%%%%%%%%%%%%
\item \textbf{Did the interaction earn its place?}

\begin{parts}
\item Compare the two models on adjusted $R^2$ and on AIC. What does each comparison say?
\answerspace{4em}

\item You are told the plain $R^2$ also rose when the interaction was added. Why is that fact
worthless on its own?
\answerspace{3.5em}

\item Give a third piece of evidence, already computed in Question 1, that supports keeping the term.
\answerspace{3em}

\item Suppose instead the interaction's interval had been $[-0.10,\; 0.42]$ and AIC had risen by
$1.4$. What would you report, and would you say the slopes are equal across types?
\answerspace{4em}
\end{parts}

%%%%%%%%%%%%%%%%%%%%%%%%%%%%%%%%%%%%%%
\item \textbf{Two intervals at one point.} For a \textbf{Branch} branch open $50$ hours a week, the
fitted value is $68.552$ with $\text{SE}_{\text{fit}} = 1.2845$.

\begin{parts}
\item Compute the 95\% confidence interval for the mean, and its width.
\answerspace{4em}

\item Compute $\sqrt{\hat\sigma^2 + \text{SE}_{\text{fit}}^2}$, then the 95\% prediction interval
for a single new branch, and its width.
\answerspace{4.5em}

\item The PI is about seven times as wide as the CI. Which of the two terms under the square root
is responsible, and why does it dominate here?
\answerspace{4em}

\item The city funds a census of every branch in the country, so $n$ becomes enormous. What happens
to each interval's width? Give the number the PI half-width cannot go below, using
$t^{\ast}\hat\sigma$.
\answerspace{4.5em}

\item The director asks: ``If I extend one branch to 50 hours, how many visits will \emph{it} get?''
Which interval do you hand her?
\answerspace{3em}
\end{parts}

%%%%%%%%%%%%%%%%%%%%%%%%%%%%%%%%%%%%%%
\item \textbf{Four criteria, four answers.} A separate study of $n = 30$ branches fits polynomials
of increasing degree. The true relationship is quadratic. Cross-validation is 6-fold.

\begin{center}
\renewcommand{\arraystretch}{1.25}
\begin{tabular}{ccccc}
\toprule
\textbf{Degree} & \textbf{$R^2$} & \textbf{Adj.\ $R^2$} & \textbf{AIC} & \textbf{CV RMSE} \\
\midrule
1 & 0.4447 & 0.4249 & 145.54 & 2.686 \\
2 & 0.5781 & 0.5469 & 139.30 & 2.471 \\
3 & 0.6131 & 0.5685 & 138.70 & 2.539 \\
5 & 0.6305 & 0.5535 & 141.32 & 2.855 \\
8 & 0.6319 & 0.4916 & 147.21 & 4.673 \\
9 & 0.7080 & 0.5766 & 142.26 & 4.117 \\
10 & 0.7145 & 0.5643 & 143.58 & 4.008 \\
\bottomrule
\end{tabular}
\end{center}

\begin{parts}
\item Fill in which degree each criterion selects.

\begin{center}
\renewcommand{\arraystretch}{1.6}
\begin{tabular}{lc}
\toprule
\textbf{Criterion} & \textbf{Degree chosen} \\
\midrule
$R^2$ & \blankline{1.1in} \\
Adjusted $R^2$ & \blankline{1.1in} \\
AIC & \blankline{1.1in} \\
CV RMSE & \blankline{1.1in} \\
\bottomrule
\end{tabular}
\end{center}

\item Only one criterion recovers the truth. Which, and what does it do that the other three
do not?
\answerspace{4em}

\item The in-sample residual standard error is $2.275$ at degree $2$ and $2.231$ at degree $10$, so
the degree-10 model fits slightly \emph{better}. Its CV RMSE is $4.008$ against $2.471$. Explain
how both statements can hold.
\answerspace{4.5em}

\item AIC picked a lower degree than adjusted $R^2$ did, but still not the truth. What does that
tell you about the difference between charging for parameters and holding data back?
\answerspace{4em}
\end{parts}

%%%%%%%%%%%%%%%%%%%%%%%%%%%%%%%%%%%%%%
\item \textbf{Find the bugs.} A classmate answers the director's question and then picks a
polynomial degree. The code runs, prints plausible numbers, and is wrong in \textbf{three} separate
ways.

\begin{lstlisting}[language=R]
fit <- lm(visits ~ hours * type, data = st)

predict(fit, newdata = data.frame(hours = 50, type = "Branch"),
        interval = "confidence")

for (dg in 1:10) {
  f <- lm(y ~ poly(x, dg), data = dd)
  cat(dg, summary(f)$r.squared, "\n")
}
\end{lstlisting}

\begin{parts}
\item Bug 1 is \texttt{interval = "confidence"}, given what the director asked in Question 3(E).
What should it be, and roughly how much wider is the right answer?
\answerspace{4em}

\item Bug 2 is choosing the degree by \texttt{r.squared}. Which degree will this loop always pick,
and why is that guaranteed before you run it?
\answerspace{4em}

\item Bug 3 is what the loop never does. Name it, and say what the loop would have to look like
instead.
\answerspace{4.5em}
\end{parts}

\end{questions}

\end{document}
